\documentclass[12pt]{article}
\usepackage[utf8]{inputenc}
\usepackage[T1]{fontenc}
\usepackage[margin=1in]{geometry}
\usepackage{setspace}
\usepackage{xcolor}
\usepackage[colorlinks=true, urlcolor=blue!60!black, linkcolor=blue!60!black]{hyperref}
\setlength{\parindent}{0em}
\setlength{\parskip}{0.45em}
\renewcommand{\baselinestretch}{1.0}
\begin{document}
\begin{center}
{\Large \textit{United States v. Darby}, 312 U.S. 100 (1941)\par}
\vspace{0.3em}
{\normalsize Case Brief --- One-Page Edition\par}
\end{center}

\textbf{Memory Jogger.} The Fair Labor Standards Act case that overruled \textit{Hammer v.\ Dagenhart} (\href{https://supreme.justia.com/cases/federal/us/312/100/}{312 U.S. 100}).

\textbf{Facts.} The FLSA set wage and hour standards for production of goods for interstate commerce and barred shipping goods made in violation. Georgia lumber manufacturer Fred Darby, who shipped lumber out of state, was indicted for underpaying and overworking his workers and failing to keep records.

\textbf{Procedural History.} The district court quashed the indictment under \textit{Hammer}; direct appeal.

\textbf{Votes.} Reversed, 8--0. Stone wrote for a unanimous Court.

\textbf{Holding.} Congress may bar interstate shipment of goods made under substandard conditions, whatever its motive, and directly regulate intrastate production that substantially affects commerce. The Tenth Amendment adds no independent limit.

\textbf{Reasoning.} The commerce power is complete and includes prohibition (\textit{Gibbons}); regulating production is an appropriate means under \textit{McCulloch} and addresses unfair interstate competition; and the Tenth Amendment is merely declaratory.

\textbf{Examples.} \emph{Same-Side:} federal safety rules for goods shipped interstate. \emph{Opposite-Side:} gun possession near schools (\textit{Lopez}). \emph{Fence-Sitter:} a purely local shop.

\textbf{Critique.} Originalists say \textit{Darby} created a federal police power; progressives credit it with recognizing a national economy.

\textbf{Key Quotation.} ``The amendment states but a truism that all is retained which has not been surrendered. There is nothing in the history of its adoption to suggest that it was more than declaratory of the relationship between the national and state governments as it had been established by the Constitution before the amendment, or that its purpose was other than to allay fears that the new national government might seek to exercise powers not granted, and that the states might not be able to exercise fully their reserved powers.'' --- Justice Stone, Opinion of the Court, 312 U.S. at 124.

\end{document}
